% Target ~350 words.

\subsection{Datasets}
\textbf{TORGO}~\cite{Rudzicz2012TORGO}: 15 native-Canadian-English
speakers (8 dysarthric, 7 controls) with 7{,}826 headset-mic utterances
across 794 unique prompts. Severity labels are from the
TORGO documentation: 5 severe (level 2), 3 mild--moderate (level 1), 7
controls (level 0). All audio resampled to 16~kHz and normalised to
$-20$~dBFS.
\textbf{UASpeech}~\cite{Kim2008UASpeech}: 15 dysarthric speakers across
isolated words, with three-listener orthographic intelligibility ratings
($r$).

\subsection{Models and features}
The headline metric uses microsoft/wavlm-large in evaluation mode, FP16
precision, with all parameters frozen and \texttt{output\_hidden\_states=True}.
For 9-layer extraction, GPU memory is the binding constraint at FP16 ($\sim$3~GB
per utterance peak). Pre-extracted embeddings are stored as
\texttt{(n\_layers, T\_frames, 1024)} \texttt{.npy} arrays under
\texttt{data/embeddings/wavlm\_large/<speaker>/}.

\subsection{Baselines}
All trained baselines are tuned fairly: the regression head \emph{and} its
hyperparameters are selected by nested group-CV (by speaker) inside each
leave-one-speaker-out fold, with a \texttt{StandardScaler} fit on the
inner-train split only (\texttt{results/finetuned\_ssl\_baseline.json}).
\textbf{MFCC + RBF-SVR}: 58-d MFCC/acoustic features mean-pooled per
utterance into an RBF Support Vector Regressor ($C$ and $\gamma$ tuned per
fold). \textbf{WavLM pool + Ridge} and \textbf{WavLM pool + LinearSVR}:
WavLM Layer-9 mean-pooled embedding into a Ridge ($\alpha$ tuned per fold)
or LinearSVR ($C$ tuned per fold) head; these are \emph{trained upper
references} for the frozen-feature setting. \textbf{MFCC trajectory DTW}:
same DTW pipeline as ATD but with MFCCs replacing SSL features, included as
the most informative ablation. We do not fine-tune the backbone: full
fine-tuning is degenerate at $n=15$ speakers, so the frozen-backbone +
trained-head probe is the fair strong baseline (following the spirit of
\cite{Hernandez2020SSLDysarthria,Schu2023SSLDys}).

\subsection{Evaluation protocol}
Speaker-level Spearman $\rho(\overline{\mathrm{ATD}}, \mathrm{severity})$
over all 15 TORGO speakers is the primary endpoint, with $p$-values from
\texttt{scipy.stats.spearmanr} and 10{,}000-sample bootstrap 95\%~CIs.
RMSE is reported on speaker-level mean predictions versus categorical
severity, on the $0$--$1$ normalised severity scale. We additionally
report a \emph{dysarthric-subgroup} Spearman over the 8 dysarthric
speakers only (controls cannot be held out as they are reference
builders); we stress this is a single $n = 8$ rank correlation, not a
leave-one-speaker-out cross-validation score, and we report its
non-significance ($p = 0.33$) and pairwise rank concordance ($\approx
73\%$) as directional indicators only. For the trained baseline heads
(MFCC + RBF-SVR, WavLM pool + Ridge/LinearSVR) the head and its
hyperparameters are refit on the remaining speakers in each
leave-one-speaker-out fold; ATD is fold-free (zero-shot), so its LOSO score
equals its full-set score.

\subsection{Accent-versus-dysarthria specificity}
To address the standing concern that any deviation-based metric may
merely measure pronunciation atypicality rather than dysarthric
pathology, we run the same ATD pipeline on two neurologically healthy
non-canonical-accent corpora: L2-ARCTIC (24 second-language English
speakers, 6 L1 backgrounds, 3{,}599 phone-annotated utterances) and
VCTK Corpus 0.92 (110 native English speakers, 13 regional accents,
restricted to the 50 sentence IDs read by $\geq$100 speakers, yielding
5{,}437 utterances). Within each corpus we build per-sentence LOCO
references by linear-interpolating all speakers' Layer-9 trajectories
to the median sentence length, averaging, and subtracting the target
speaker's contribution. DTW uses the same banded algorithm
(window~20). The specificity hypothesis is pre-registered: (i)
within L2-ARCTIC, Spearman $\rho(\text{mean ATD},\text{phone-error
rate})$ should not differ from zero; (ii) within VCTK, Kruskal--Wallis
across accent groups should not be significant; (iii) on the pooled
4-corpus dataset, the dysarthric-vs-healthy contrast on
within-corpus-z-normalised ATD should be highly significant.

\subsection{Confound and sensitivity analyses}
We check that ATD is not a proxy for speaking rate by reporting the
Pearson correlation between ATD and per-utterance duration.
Per-speaker leverage is reported via leave-one-out on each of the 15
speakers and via Cook's-distance-style sensitivity in
the supplementary. We additionally pre-register a controlled-perturbation
study on the seven control speakers (\S~5.4 and Appendix~B); its
pre-registered articulatory-specificity hypothesis is \emph{falsified}
(pink noise outranks formant shift at both layers, triggering locked
falsification criterion~\#2), and we report it as such in the Limitations
section.
